% ═══════════════════════════════════════════════════════════════
%  chapters/ch13-scope-render.tex
%  Отрисовка осциллограммы через Cairo
% ═══════════════════════════════════════════════════════════════

\chapter{Отрисовка осциллограммы}
\label{ch:scope-render}

\epigraph{<<Осциллограмма~--- это не просто линия. Это \emph{язык},
на котором сигнал говорит с инженером>>.}{}

\section{Введение}
\label{sec:render-intro}

В \texttt{ch12} мы сделали \emph{потоки}, \emph{мьютексы},
\emph{кольцо} \emph{кадров}. Данные \emph{текут}. Но
\emph{пользователь} \emph{их} \emph{не} \emph{видит}. Пора
\emph{нарисовать} \emph{осциллограмму}.

\textbf{Инструмент}~--- \textbf{Cairo}. \emph{Тот} \emph{же},
что \emph{в} \emph{первой} \emph{книге}. \emph{Но} \emph{задача}
\emph{другая}: не \emph{график} \emph{фоторезистора}, а
\emph{осциллограмма} \emph{сигнала}.

\begin{ideabox}
\textbf{Осциллограмма~--- это \emph{окно} в мир сигнала.}
Она \emph{показывает} \emph{форму}, \emph{амплитуду},
\emph{частоту}. \emph{Хорошая} \emph{осциллограмма}~---
\emph{читается} \emph{мгновенно}. \emph{Плохая}~---
\emph{требует} \emph{расшифровки}.
\end{ideabox}

\subsection{Что вы узнаете}

\begin{itemize}
    \item Как \emph{перевести} данные в \emph{координаты} \emph{экрана}.
    \item Как \emph{нарисовать} \emph{сетку} и \emph{оси}.
    \item Как \emph{построить} \emph{линию} \emph{графика}.
    \item Как \emph{добавить} \emph{статистику} и \emph{легенду}.
    \item Как \emph{сделать} \emph{кнопку} \emph{<<Очистить>>}
          \emph{на} \emph{canvas}.
\end{itemize}

\section{Координаты: данные в пиксели}
\label{sec:render-coords}

\textbf{Первая} \emph{задача}~--- \emph{преобразовать}
\emph{данные} \emph{в} \emph{координаты} \emph{экрана}.

\textbf{У} \emph{нас} \emph{есть}:
\begin{itemize}
    \item \textbf{Значения} \emph{АЦП} (\num{0}--\num{4095}).
    \item \textbf{Номер} \emph{точки} ($0$ -- $N-1$).
\end{itemize}

\textbf{Надо} \emph{получить}:
\begin{itemize}
    \item \textbf{$X$}~--- \emph{от} \texttt{left\_margin}
          \emph{до} \texttt{width - right\_margin}.
    \item \textbf{$Y$}~--- \emph{от} \texttt{height - bottom\_margin}
          \emph{до} \texttt{top\_margin}.
\end{itemize}


\subsection{Формулы}

\textbf{Для} \emph{$X$}:
\[
x = \text{left\_margin} + i \cdot \frac{\text{graph\_w}}{N - 1}
\]
где \(i\)~--- \emph{номер} \emph{точки}, \(N\)~--- \emph{количество}
\emph{точек}.

\textbf{Для} \emph{$Y$}:
\[
y = \text{height} - \text{bottom\_margin}
    - \frac{v}{\text{ADC\_MAX}} \cdot \text{graph\_h}
\]
где \(v\)~--- \emph{значение} \emph{АЦП}.

\begin{keybox}
\textbf{Cairo} \emph{считает} \emph{$Y$} \emph{сверху}
\emph{вниз}. \emph{Поэтому} \emph{мы} \emph{вычитаем} \emph{из}
\texttt{height - bottom\_margin}. \emph{Иначе} \emph{график}
\emph{будет} \emph{перевёрнут}.
\end{keybox}

\section{Сетка и оси}
\label{sec:render-grid}

\textbf{Сетка}~--- \emph{каркас} \emph{графика}. \emph{Без}
\emph{неё}~--- \emph{облако} \emph{точек}. \emph{С} \emph{ней}~---
\emph{понятно}, \emph{где} \emph{какое} \emph{значение}.

\begin{listing}[htbp]
\begin{minted}[fontsize=\footnotesize, frame=none, bgcolor=gray!5]{c}
// Сетка по Y (8 линий)
for (int i = 0; i <= 8; i++) {
    double y = height - bottom_margin
             - (i * graph_h / 8.0);
    int value = i * ADC_MAX / 8;

    cairo_set_source_rgb(cr, 0.9, 0.9, 0.9);
    cairo_set_line_width(cr, 1);
    cairo_move_to(cr, left_margin, y);
    cairo_line_to(cr, width - right_margin, y);
    cairo_stroke(cr);

    // Метка
    char label[16];
    snprintf(label, sizeof(label), "%d", value);
    cairo_set_source_rgb(cr, 0, 0, 0);
    cairo_move_to(cr, left_margin - 38, y + 4);
    cairo_show_text(cr, label);
}
\end{minted}
\caption{Сетка по Y}
\label{lst:render-grid-y}
\end{listing}

\textbf{Что} \emph{здесь} \emph{важно}:
\begin{itemize}
    \item \textbf{8} \emph{линий}~--- \emph{универсальный}
          \emph{выбор}. \emph{Не} \emph{слишком} \emph{часто},
          \emph{не} \emph{слишком} \emph{редко}.
    \item \textbf{Светло-серый} \emph{цвет}~--- \emph{сетка}
          \emph{не} \emph{должна} \emph{затмевать} \emph{сигнал}.
    \item \textbf{Метки} \emph{слева}~--- \emph{значения}
          \emph{АЦП}.
\end{itemize}

\section{Линия графика}
\label{sec:render-line}

\textbf{Линия}~--- \emph{главный} \emph{элемент}. \emph{Она}
\emph{соединяет} \emph{соседние} \emph{точки} \emph{и}
\emph{показывает} \emph{форму} \emph{сигнала}.

\begin{listing}[htbp]
\begin{minted}[fontsize=\footnotesize, frame=none, bgcolor=gray!5]{c}
cairo_set_source_rgb(cr, 0.2, 1.0, 0.2);  // ярко-зелёный
cairo_set_line_width(cr, 1.5);

// Копируем данные под мьютексом
OscFrame frame;
if (ring_pop(&app->ring, &frame)) {
    int n = frame.nsamples;
    double x_step = (n > 1) ? (double)graph_w / (n - 1) : 0;

    for (int i = 0; i < n; i++) {
        double x = left_margin + i * x_step;
        double y = height - bottom_margin
                 - (frame.samples[i] / ADC_MAX) * graph_h;

        if (i == 0) cairo_move_to(cr, x, y);
        else        cairo_line_to(cr, x, y);
    }
    cairo_stroke(cr);
}
\end{minted}
\caption{Линия графика}
\label{lst:render-line}
\end{listing}

\begin{warnbox}
\textbf{Данные} \emph{берём} \emph{под} \emph{мьютексом}.
\emph{Иначе}~--- \emph{гонка} \emph{данных}. \emph{График}
\emph{может} \emph{<<сломаться>>} \emph{посередине}.
\end{warnbox}

\section{Статистика и легенда}
\label{sec:render-stats}

\textbf{Статистика}~--- \emph{четыре} \emph{числа}, \emph{которые}
\emph{говорят} \emph{о} \emph{сигнале} \emph{больше}, \emph{чем}
\emph{тысяча} \emph{пикселей}:
\begin{itemize}
    \item \textbf{Текущее}~--- \emph{последнее} \emph{значение}.
    \item \textbf{Минимум}~--- \emph{самое} \emph{низкое}.
    \item \textbf{Максимум}~--- \emph{самое} \emph{высокое}.
    \item \textbf{Среднее}~--- \emph{типичное}.
\end{itemize}

\begin{listing}[htbp]
\begin{minted}[fontsize=\footnotesize, frame=none, bgcolor=gray!5]{c}
PlotStats s = compute_stats(&frame);
if (s.valid) {
    char stats[128];
    snprintf(stats, sizeof(stats),
             "Тек: %d  Мин: %d  Макс: %d  Сред: %d",
             s.last, s.min, s.max, s.avg);

    cairo_set_source_rgb(cr, 0, 0.5, 0);
    cairo_set_font_size(cr, 10);
    cairo_move_to(cr, left_margin, top_margin - 10);
    cairo_show_text(cr, stats);
}
\end{minted}
\caption{Статистика}
\label{lst:render-stats}
\end{listing}

\begin{ideabox}
\textbf{Легенда}~--- \emph{ключ} \emph{к} \emph{графику}. \emph{Она}
\emph{объясняет}, \emph{что} \emph{означает} \emph{зелёная}
\emph{линия}. \emph{Без} \emph{неё}~--- \emph{<<чёрный} \emph{ящик>>}.
\end{ideabox}

\section{Кнопка <<Очистить>> на canvas}
\label{sec:render-button}

\textbf{Кнопка} \emph{<<Очистить>>}~--- \emph{не} \emph{GTK-виджет}.
\emph{Она} \emph{нарисована} \emph{прямо} \emph{на} \emph{canvas}.
\emph{Клик} \emph{обрабатывается} \emph{вручную}.

\begin{listing}[htbp]
\begin{minted}[fontsize=\footnotesize, frame=none, bgcolor=gray!5]{c}
// Отрисовка
double btn_x = width / 2.0 - 80;
double btn_y = 5;
double btn_w = 160;
double btn_h = 28;

cairo_set_source_rgb(cr, 0.8, 0.8, 0.8);
cairo_rectangle(cr, btn_x, btn_y, btn_w, btn_h);
cairo_fill(cr);

cairo_set_source_rgb(cr, 0, 0, 0);
cairo_set_font_size(cr, 12);
cairo_move_to(cr, btn_x + 25, btn_y + 19);
cairo_show_text(cr, "Очистить график");

// Обработка клика
gboolean on_graph_button_press(GtkWidget *w,
                               GdkEventButton *e,
                               gpointer data)
{
    AppData *app = data;
    int width = gtk_widget_get_allocated_width(w);
    double btn_x = width / 2.0 - 80;
    double btn_y = 5;
    double btn_w = 160;
    double btn_h = 28;

    if (e->x >= btn_x && e->x <= btn_x + btn_w &&
        e->y >= btn_y && e->y <= btn_y + btn_h) {
        ring_clear(&app->ring);
        gtk_widget_queue_draw(app->drawing_area);
        return TRUE;
    }
    return FALSE;
}
\end{minted}
\caption{Кнопка <<Очистить>>: отрисовка и обработка}
\label{lst:render-button}
\end{listing}



\begin{figure}[htbp]
\centering
\adjustbox{max width=\textwidth}{%
\begin{tikzpicture}[
    x=1cm, y=1cm
]

% ─── Координатная сетка (для отладки) ────────────
\draw[step=1cm, gray!20, very thin] (-1, -1) grid (13, 7);

% ─── Внешний прямоугольник — виджет ──────────────
\draw[thick, fill=white] (0, 0) rectangle (12, 6);

% ─── Внутренний прямоугольник — область графика ──
\draw[thick, fill=clrLogicData!10] (2, 1) rectangle (10, 5);

% Название внутри
\node[font=\small\bfseries] at (6, 3) {график};

% ─── Подписи отступов (снаружи) ──────────────────
\node[font=\footnotesize\ttfamily, gray!70!black, anchor=south]
    at (6, 5.1) {top\_margin};

\node[font=\footnotesize\ttfamily, gray!70!black, anchor=north]
    at (6, 0.9) {bottom\_margin};

\node[font=\footnotesize\ttfamily, gray!70!black, rotate=90, anchor=south]
    at (1.9, 3) {left\_margin};

\node[font=\footnotesize\ttfamily, gray!70!black, rotate=-90, anchor=north]
    at (10.6, 3) {right\_margin};

% Метки осей (для отладки)
\foreach \x in {-1,...,13} {
    \node[font=\tiny, gray] at (\x, -1.3) {\x};
}
\foreach \y in {-1,...,7} {
    \node[font=\tiny, gray] at (-1.3, \y) {\y};
}

\end{tikzpicture}
}% end adjustbox
\caption{Схема отступов в графике.}
\label{fig:render-margins}
\end{figure}



\begin{warnbox}
\textbf{Координаты} \emph{кнопки} \emph{дублируются}~---
\emph{в} \emph{отрисовке} \emph{и} \emph{в} \emph{обработке}.
\emph{Если} \emph{поменять} \emph{в} \emph{одном} \emph{месте}
\emph{и} \emph{забыть} \emph{в} \emph{другом}~--- \emph{кнопка}
\emph{<<перестанет} \emph{нажиматься>>}. \emph{Вынесите}
\emph{их} \emph{в} \emph{константы}.
\end{warnbox}

\section{Уроки главы}
\label{sec:render-lessons}

\begin{enumerate}
    \item \textbf{Координаты~--- первое.} Без \emph{преобразования}
          \emph{данных} \emph{в} \emph{пиксели}~--- \emph{ничего}
          \emph{не} \emph{нарисуешь}.
    \item \textbf{Cairo} \emph{считает} \emph{$Y$} \emph{сверху}
          \emph{вниз}. \emph{Не} \emph{забудьте} \emph{вычесть}.
    \item \textbf{Сетка~--- 8} \emph{линий}. \emph{Универсальный}
          \emph{выбор}.
    \item \textbf{Данные} \emph{берём} \emph{под} \emph{мьютексом}.
          \emph{Иначе}~--- \emph{гонка}.
    \item \textbf{Статистика~--- 4} \emph{числа}. \emph{Больше},
          \emph{чем} \emph{1000} \emph{пикселей}.
    \item \textbf{Кнопка} \emph{на} \emph{canvas}~--- \emph{два}
          \emph{места}: \emph{отрисовка} \emph{и} \emph{клик}.
          \emph{Координаты} \emph{должны} \emph{совпадать}.
\end{enumerate}

\section{Чек-лист грабель}
\label{sec:render-pitfalls}

\begin{table}[htbp]
\centering
\caption{Грабли отрисовки}
\label{tab:render-pitfalls}
\small
\begin{tabularx}{\textwidth}{@{}c L L@{}}
\toprule
\textbf{№} & \textbf{Грабля} & \textbf{Решение} \\
\midrule
1 & График перевёрнут & \texttt{height - bottom\_margin} \\
2 & Нет сетки & 8 линий по Y \\
3 & Гонка данных & \texttt{g\_mutex\_lock()} \\
4 & Кнопка не нажимается & Координаты в константы \\
5 & Забыли \texttt{cairo\_stroke} & Линия не рисуется \\
6 & Метки налезают & Сдвиг \texttt{left\_margin - 38} \\
7 & Нет \texttt{queue\_draw} & График не обновляется \\
8 & Нет проверки \texttt{valid} & Деление на 0 \\
\bottomrule
\end{tabularx}
\end{table}

\section{Упражнения}
\label{sec:render-exercises}

\begin{exercisebox}
\textbf{Упражнение 1.} Нарисуйте \emph{вертикальную}
\emph{линию}~--- \emph{маркер} \emph{максимума}. \emph{Цвет}~---
\emph{красный}, \emph{толщина}~--- \emph{2}.
\end{exercisebox}

\begin{exercisebox}
\textbf{Упражнение 2.} Добавьте \emph{масштабирование}:
\emph{колёсико} \emph{мыши} \emph{меняет} \emph{масштаб} \emph{по}
\emph{$X$}.
\end{exercisebox}

\needspace{5\baselineskip}
\begin{exercisebox}
\textbf{Упражнение 3.} Реализуйте \emph{курсор}: \emph{при}
\emph{наведении} \emph{показывает} \emph{значение} \emph{в}
\emph{точке}.
\end{exercisebox}

\begin{exercisebox}
\textbf{Упражнение 4.} Сохраните \emph{осциллограмму} в
\emph{PNG}. \emph{Какие} \emph{API} \emph{нужны}?
\end{exercisebox}

\section{Проверьте сами}
\label{sec:render-check}

\begin{enumerate}
    \item Соберите \emph{ПК-приложение}.
    \item Подключите \emph{Nucleo}.
    \item Запустите \emph{поток}.
    \item \emph{Проверьте}, что \emph{график} \emph{рисуется}.
    \item \emph{Нажмите} <<Очистить>>~--- \emph{график} \emph{пуст}.
    \item \emph{Сделайте} \emph{скриншот}~--- \emph{пришлите}
          \emph{в} \emph{книгу}.
\end{enumerate}

\begin{ideabox}
\textbf{Если получилось~--- вы поняли Cairo.} Это
\emph{главное}. Дальше~--- \emph{интерфейс} \emph{генератора},
\emph{жизненный} \emph{цикл}, \emph{отладка}. \emph{Всё}
\emph{впереди}.
\end{ideabox}


